\documentclass[12pt, a4paper]{article}

% --- Ngôn ngữ & Phông chữ ---
\usepackage[utf8]{inputenc}
\usepackage[utf8]{vietnam} % Chuẩn hiển thị tiếng Việt cho pdflatex

% --- Cấu trúc trang & Lề ---
\usepackage{geometry}
\usepackage{fancyhdr}
\usepackage{lastpage}
\usepackage{setspace}
\usepackage{parskip}

% --- Bảng biểu ---
\usepackage{array}
\usepackage{tabularx}
\usepackage{longtable}
\usepackage{booktabs}
\usepackage{multirow}
\usepackage{multicol}
\usepackage{diagbox}
\usepackage{rotating}

% --- Toán học & Ký hiệu ---
\usepackage{amsmath, amsfonts, amssymb, amsthm}
\usepackage{wasysym}
\usepackage[makeroom]{cancel}

% --- Đồ họa & Màu sắc ---
\usepackage{xcolor}
\usepackage{graphicx}
\usepackage{caption}
\usepackage{subcaption}
\usepackage{float}

% --- TikZ & Vẽ đồ thị ---
\usepackage{tikz}
\usepackage{pgfplots}
\usepackage{pgfgantt}
\usepackage{dirtree}
\pgfplotsset{compat=1.18}
\usetikzlibrary{
    3d, calc, positioning, arrows.meta, 
    shapes, shapes.geometric, automata, petri
}

% --- Khung & Mã nguồn ---
\usepackage{alltt}
\usepackage{listings}
\usepackage{tcolorbox}
\usepackage[framemethod=tikz]{mdframed}
\usepackage[lined,boxed,commentsnumbered]{algorithm2e}

% --- Tiện ích khác ---
\usepackage{enumitem}
\usepackage{titling}
\usepackage{hyperref}
\usepackage[backend=biber, style=ieee]{biblatex}
\addbibresource{references.bib}

% --- Nạp cấu hình từ file Setup.tex ---
% ===== Cấu hình căn lề A4 chuẩn luận văn =====
\geometry{
    top=2.5cm,
    bottom=3.0cm,
    left=2.0cm,
    right=2.0cm
}
\hypersetup{urlcolor=blue,linkcolor=black,citecolor=black,colorlinks=true} 
%\usepackage{pstcol} % PSTricks with the standard color package

\usepackage[normalem]{ulem}

\newtheorem{theorem}{{\bf Định lý}}

\newtheorem{property}{{\bf Tính chất}}
\newtheorem{proposition}{{\bf Mệnh đề}}
\newtheorem{corollary}[proposition]{{\bf Hệ quả}}
\newtheorem{lemma}[proposition]{{\bf Bổ đề}}
\theoremstyle{definition}
\newtheorem{exer}{Bài toán}


\setcounter{secnumdepth}{4}
\setcounter{tocdepth}{3}
\makeatletter
\newcounter {subsubsubsection}[subsubsection]
\renewcommand\thesubsubsubsection{\thesubsubsection .\@alph\c@subsubsubsection}
\newcommand\subsubsubsection{\@startsection{subsubsubsection}{4}{\z@}%
                                     {-3.25ex\@plus -1ex \@minus -.2ex}%
                                     {1.5ex \@plus .2ex}%
                                     {\normalfont\normalsize\bfseries}}
\newcommand*\l@subsubsubsection{\@dottedtocline{3}{10.0em}{4.1em}}
\newcommand*{\subsubsubsectionmark}[1]{}
\makeatother

\everymath{\color{black}}%make in-line maths symbols blue to read/check easily

\sloppy
\captionsetup[figure]{labelfont={small,bf},textfont={small,it},belowskip=-1pt,aboveskip=-9pt}
%space remove between caption, figure, and text
\captionsetup[table]{labelfont={small,bf},textfont={small,it},belowskip=-1pt,aboveskip=7pt}


\setlength{\floatsep}{5pt plus 2pt minus 2pt}
\setlength{\textfloatsep}{5pt plus 2pt minus 2pt}
\setlength{\intextsep}{10pt plus 2pt minus 2pt}


\lstset{
language=Matlab,
    aboveskip=6mm,
    belowskip=-6mm,
    backgroundcolor=\color{white},   
    commentstyle=\color{codegreen},
    keywordstyle=\color{blue},
    numberstyle=\tiny\color{codegray},
    stringstyle=\color{codepurple},
    basicstyle=\ttfamily,
    breakatwhitespace=false,         
    breaklines=true,                 
    captionpos=b,                    
    keepspaces=true,                 
    numbers=left,                    
    numbersep=5pt,                  
    showspaces=false,                
    showstringspaces=false,
    showtabs=false,                  
    tabsize=2,
    frame=tb,
    framesep=8pt,
    rulecolor=\color{black},
    title=\lstname,
}

%New colors defined below
\definecolor{codegreen}{rgb}{0,0.6,0}
\definecolor{codegray}{rgb}{0.5,0.5,0.5}
\definecolor{codepurple}{rgb}{0.58,0,0.82}
\definecolor{backcolour}{rgb}{0.95,0.95,0.92}

%Code listing style named "mystyle"
\lstdefinestyle{mystyle}{
  backgroundcolor=\color{backcolour}, 
  commentstyle=\color{codegreen},
  keywordstyle=\color{magenta},
  numberstyle=\tiny\color{codegray},
  stringstyle=\color{codepurple},
  basicstyle=\ttfamily\footnotesize,
  breakatwhitespace=false,         
  breaklines=true,                 
  captionpos=b,                    
  keepspaces=true,                 
  numbers=left,                    
  numbersep=5pt,                  
  showspaces=false,                
  showstringspaces=false,
  showtabs=false,                  
  tabsize=2
}

%"mystyle" code listing set
\lstset{style=mystyle}
\lstset{
        literate=%
% Vần a
	{á}{{\'a}}1
	{à}{{\`a}}1
	{ạ}{{\d a}}1
	{ả}{{\h a}}1
	{ã}{{\~ a}}1
	%
	{Á}{{\'A}}1
	{À}{{\`A}}1
	{Ạ}{{\d A}}1
	{Ả}{{\h A}}1
	{Ã}{{\~ A}}1
%
% Vần ă
	{ă}{{\u a}}1
	{ắ}{{\'\abreve }}1
	{ằ}{{\`\abreve }}1
	{ặ}{{\d \abreve }}1
	{ẳ}{{\h \abreve }}1
	{ẵ}{{\~\abreve }}1
	%
	{Ă}{{\u A}}1
	{Ắ}{{\'\ABREVE }}1
	{Ằ}{{\`\ABREVE }}1
	{Ặ}{{\d \ABREVE }}1
	{Ẳ}{{\h \ABREVE }}1
	{Ẵ}{{\~\ABREVE }}1
%
% Vần â
	{â}{{\^ a}}1
	{ấ}{{\'\acircumflex }}1
	{ầ}{{\`\acircumflex }}1
	{ậ}{{\d \acircumflex }}1
	{ẩ}{{\h \acircumflex }}1
	{ẫ}{{\~\acircumflex }}1
	 %
	{Â}{{\^ A}}1
	{Ấ}{{\'\ACIRCUMFLEX }}1
	{Ầ}{{\`\ACIRCUMFLEX }}1
	{Ậ}{{\d \ACIRCUMFLEX }}1
	{Ẩ}{{\h \ACIRCUMFLEX }}1
	{Ẫ}{{\~\ACIRCUMFLEX }}1
%
% Vần đ
	{đ}{{\dj }}1
	{Đ}{{\DJ }}1
%
% Vần e
	{é}{{\'e}}1
	{è}{{\`e}}1
	{ẹ}{{\d e}}1
	{ẻ}{{\h e}}1
	{ẽ}{{\~ e}}1
	%
	{É}{{\'E}}1
	{È}{{\`E}}1
	{Ẹ}{{\d E}}1
	{Ẻ}{{\h E}}1
	{Ẽ}{{\~ E}}1
%
% Vần ê
	{ê}{{\^e}}1
	{ế}{{\'\ecircumflex }}1
	{ề}{{\`\ecircumflex }}1
	{ệ}{{\d \ecircumflex }}1
	{ể}{{\h \ecircumflex }}1
	{ễ}{{\~\ecircumflex }}1
	%
	{Ê}{{\^E}}1
	{Ế}{{\'\ECIRCUMFLEX }}1
	{Ề}{{\`\ECIRCUMFLEX }}1
	{Ệ}{{\d \ECIRCUMFLEX }}1
	{Ể}{{\h \ECIRCUMFLEX }}1
	{Ễ}{{\~\ECIRCUMFLEX }}1
%
% Vần i
	{í}{{\'i}}1
	{ì}{{\`\i }}1
	{ị}{{\d i}}1
	{ỉ}{{\h i}}1
	{ĩ}{{\~\i }}1
	%
	{Í}{{\'I}}1
	{Ì}{{\`I}}1
	{Ị}{{\d I}}1
	{Ỉ}{{\h I}}1
	{Ĩ}{{\~I}}1
%
% Vần o
	{ó}{{\'o}}1
	{ò}{{\`o}}1
	{ọ}{{\d o}}1
	{ỏ}{{\h o}}1
	{õ}{{\~o}}1
	%
	{Ó}{{\'O}}1
	{Ò}{{\`O}}1
	{Ọ}{{\d O}}1
	{Ỏ}{{\h O}}1
	{Õ}{{\~O}}1
%
% Vần ô
	{ô}{{\^o}}1
	{ố}{{\'\ocircumflex }}1
	{ồ}{{\`\ocircumflex }}1
	{ộ}{{\d \ocircumflex }}1
	{ổ}{{\h \ocircumflex }}1
	{ỗ}{{\~\ocircumflex }}1
	%
	{Ô}{{\^O}}1
	{Ố}{{\'\OCIRCUMFLEX }}1
	{Ồ}{{\`\OCIRCUMFLEX }}1
	{Ộ}{{\d \OCIRCUMFLEX }}1
	{Ổ}{{\h \OCIRCUMFLEX }}1
	{Ỗ}{{\~\OCIRCUMFLEX }}1
%
% Vần ơ
	{ơ}{{\ohorn }}1
	{ớ}{{\'\ohorn }}1
	{ờ}{{\`\ohorn }}1
	{ợ}{{\d \ohorn }}1
	{ở}{{\h \ohorn }}1
	{ỡ}{{\~\ohorn }}1
	%
	{Ơ}{{\OHORN }}1
	{Ớ}{{\'\OHORN }}1
	{Ờ}{{\`\OHORN }}1
	{Ợ}{{\d \OHORN }}1
	{Ở}{{\h \OHORN }}1
	{Ỡ}{{\~\OHORN }}1
%
% Vần u
	{ú}{{\'u}}1
	{ù}{{\`u}}1
	{ụ}{{\d u}}1
	{ủ}{{\h u}}1
	{ũ}{{\~u}}1
	%
	{Ú}{{\'U}}1
	{Ù}{{\`U}}1
	{Ụ}{{\d U}}1
	{Ủ}{{\h U}}1
	{Ũ}{{\~U}}1
%
% Vần ư
	{ư}{{\uhorn }}1
	{ứ}{{\'\uhorn }}1
	{ừ}{{\`\uhorn }}1
	{ự}{{\d \uhorn }}1
	{ử}{{\h \uhorn }}1
	{ữ}{{\~\uhorn }}1
	%
	{Ư}{{\UHORN }}1
	{Ứ}{{\'\UHORN }}1
	{Ừ}{{\`\UHORN }}1
	{Ự}{{\d \UHORN }}1
	{Ử}{{\h \UHORN }}1
	{Ữ}{{\~\UHORN }}1
%
% Vần y
	{ý}{{\'y}}1
	{ỳ}{{\`y}}1
	{ỵ}{{\d y}}1
	{ỷ}{{\h y}}1
	{ỹ}{{\~y}}1
	%
	{Ý}{{\'Y}}1
	{Ỳ}{{\`Y}}1
	{Ỵ}{{\d Y}}1
	{Ỷ}{{\h Y}}1
	{Ỹ}{{\~Y}}1
}
 

% --- Cấu hình Header & Footer ---
\pagestyle{fancy}
\fancyhf{} 
\setlength{\headheight}{42pt}
\renewcommand{\headrulewidth}{0.4pt}
\renewcommand{\footrulewidth}{0.4pt}
\renewcommand{\footruleskip}{1.5mm}

% Header chuẩn Bách Khoa CSE
\fancyhead[L]{%
  \begin{tabular}{rl}
    \begin{picture}(25,15)(0,0)
      \put(0,-8){\includegraphics[width=10mm, height=10mm]{images/hcmut.png}}
    \end{picture}&
    \begin{tabular}{l}
      \textbf{Trường Đại Học Bách Khoa - ĐHQG TP. HCM} \\
      \textbf{Khoa Khoa Học và Kỹ Thuật Máy Tính}
    \end{tabular} 	
  \end{tabular}%
}

% Footer mặc định cho toàn bộ nội dung
\fancyfoot[L]{\small Kỹ năng Chuyên nghiệp cho Kỹ sư (CO2001) - Nhóm G2}
\fancyfoot[R]{\small Trang \thepage/\pageref{LastPage}}


% ==============================================
% BẮT ĐẦU NỘI DUNG TÀI LIỆU
% ==============================================
\begin{document}

% --- Trang bìa ---
\begin{titlepage}
    \begin{titlepage}

   \begin{tikzpicture}[remember picture,overlay,inner sep=0,outer sep=0]
     \draw[blue!70!black,line width=4pt] ([xshift=-2cm,yshift=-2cm]current page.north east) coordinate (A)--([xshift=2cm,yshift=-2cm]current page.north west) coordinate(B)--([xshift=2cm,yshift=2cm]current page.south west) coordinate (C)--([xshift=-2cm,yshift=2cm]current page.south east) coordinate(D)--cycle;
     \draw ([yshift=0.5cm,xshift=-0.5cm]A)-- ([yshift=0.5cm,xshift=0.5cm]B)--
     ([yshift=-0.5cm,xshift=0.5cm]B) --([yshift=-0.5cm,xshift=-0.5cm]B)--([yshift=0.5cm,xshift=-0.5cm]C)--([yshift=0.5cm,xshift=0.5cm]C)--([yshift=-0.5cm,xshift=0.5cm]C)-- ([yshift=-0.5cm,xshift=-0.5cm]D)--([yshift=0.5cm,xshift=-0.5cm]D)--([yshift=0.5cm,xshift=0.5cm]D)--([yshift=-0.5cm,xshift=0.5cm]A)--([yshift=-0.5cm,xshift=-0.5cm]A)--([yshift=0.5cm,xshift=-0.5cm]A);

     \draw ([yshift=-0.3cm,xshift=0.3cm]A)-- ([yshift=-0.3cm,xshift=-0.3cm]B)--
     ([yshift=0.3cm,xshift=-0.3cm]B) --([yshift=0.3cm,xshift=0.3cm]B)--
     ([yshift=-0.3cm,xshift=0.3cm]C)--([yshift=-0.3cm,xshift=-0.3cm]C)--([yshift=0.3cm,xshift=-0.3cm]C)-- ([yshift=0.3cm,xshift=0.3cm]D)--([yshift=-0.3cm,xshift=0.3cm]D)--([yshift=-0.3cm,xshift=-0.3cm]D)--([yshift=0.3cm,xshift=-0.3cm]A)--([yshift=0.3cm,xshift=0.3cm]A)--([yshift=-0.3cm,xshift=0.3cm]A);
   \end{tikzpicture}
   
\vspace*{-0.8cm}
  
\begin{center}
    {\large \textbf{ĐẠI HỌC QUỐC GIA THÀNH PHỐ HỒ CHÍ MINH}} \\
    {\large \textbf{TRƯỜNG ĐẠI HỌC BÁCH KHOA}} \\[0.2cm]
    \textbf{\large KHOA KHOA HỌC VÀ KỸ THUẬT MÁY TÍNH}
\end{center}

\vfill

\begin{center}
    \includegraphics[width=3.5cm]{images/hcmut.png}
\end{center}

\vfill

\begin{center}
    {\Huge \textbf{BÁO CÁO ĐỀ CƯƠNG}}\\[0.3cm]
    {\large \textbf{KỸ NĂNG CHUYÊN NGHIỆP CHO KỸ SƯ - CO2001}}
\end{center}

\vspace{1.5cm}

\begin{center}
    {\Large \textbf{ĐỀ TÀI:}}\\[0.3cm]
    {\Large \textbf{Kỹ năng cần thiết cho kỹ sư phần mềm}}\\[0.2cm]
    {\Large \textbf{trong kỷ nguyên AI và Machine Learning}}
\end{center}

\vspace{2cm}

\begin{center}
    {\large \textbf{NHÓM G2}}
\end{center}

\vfill

% Dòng thời gian
\centerline{\large \textit{Thành phố Hồ Chí Minh, tháng 09/2026}}
\vspace*{0.8cm} 

\end{titlepage}

\end{titlepage}

% --- Trang danh sách thành viên ---
\newpage
\thispagestyle{fancy}
\noindent\textbf{\large Giảng viên hướng dẫn: PHẠM MINH TUẤN}\\[0.5cm]
\noindent\textbf{\large Danh sách thành viên nhóm G2:}
\vspace{0.3cm}

\renewcommand{\arraystretch}{1.4}
\begin{tabularx}{\textwidth}{|p{4.2cm}|X|}
    \hline
    \textbf{Thành viên nhóm} & \textbf{Nhiệm vụ} \\
    \hline
    Đặng Minh Quân & Phân tích tài liệu; xây dựng các nhóm kỹ năng cần thiết \\ \hline
    Đặng Văn Quang & Rủi ro/hạn chế và khuyến nghị cho sinh viên, kết luận, chuẩn hoá báo cáo và slide \\ \hline
    Đào Văn Trọng Hiệp & Phân tích tài liệu; xây dựng các nhóm kỹ năng cần thiết kỹ sư phần mềm \\ \hline 
    Đinh Lê Nguyên & Chuẩn bị tài liệu, thực trạng và so sánh workflow trước 2022 và hiện tại \\ \hline
    Đào Quốc Minh & Rủi ro/hạn chế và khuyến nghị cho sinh viên, kết luận, chuẩn hoá báo cáo và slide \\ \hline
    Dương Bảo Kiệt & Chuẩn bị tài liệu, thực trạng và so sánh workflow trước 2022 và hiện tại \\ \hline
\end{tabularx}

% --- Mục lục ---
\newpage
\tableofcontents

% --- Danh mục viết tắt ---
\newpage 
\section*{DANH MỤC TỪ VIẾT TẮT}
\addcontentsline{toc}{section}{DANH MỤC TỪ VIẾT TẮT}

\begin{table}[h!]
\renewcommand{\arraystretch}{1.3}
\begin{tabular}{@{}p{2.5cm} p{7cm} p{6cm}@{}}
\textbf{Viết tắt} & \textbf{Tiếng Anh} & \textbf{Tiếng Việt} \\ \hline
\textbf{AI} & Artificial Intelligence & Trí tuệ nhân tạo \\
\textbf{GenAI} & Generative Artificial Intelligence & Trí tuệ nhân tạo tạo sinh \\
\textbf{LLM} & Large Language Model & Mô hình ngôn ngữ lớn \\
\textbf{DACUM} & Developing A CurriculUM & Phương pháp phát triển khung chương trình đào tạo \\
\textbf{SDLC} & Software Development Life Cycle & Vòng đời phát triển phần mềm \\
\textbf{PR} & Pull Request & Yêu cầu hợp nhất mã nguồn \\
\textbf{POC} & Proof of Concept & Bản thử nghiệm chứng minh khái niệm \\
\textbf{SKA} & Skills, Knowledge, and Attributes & Kỹ năng, Kiến thức và Phẩm chất \\
\textbf{CUJ} & Critical User Journey & Hành trình người dùng trọng yếu \\
\textbf{GDPR} & General Data Protection Regulation & Quy định chung về bảo vệ dữ liệu (EU) \\
\textbf{IDE} & Integrated Development Environment & Môi trường phát triển tích hợp \\ \hline
\end{tabular}
\end{table}

% --- Nội dung các chương ---
\newpage
\section{MỞ ĐẦU}
\subsection{Bối cảnh và động lực}
Ví dụ trích dẫn \cite{stackoverflow2024survey}, hay kiểu \cite{kam2025what} và báo cáo từ McKinsey \cite{mckinsey2023unleashing}  phần mềm.
v0 cho bản gốc
\subsection{Giới thiệu bài báo khoa học nền tảng}
\subsection{Mục tiêu và phạm vi của báo cáo nhóm}

Dựa trên bài báo nền tảng, báo cáo tập trung phân tích sự thay đổi trong vai trò và quy trình làm việc của kỹ sư phần mềm khi GenAI được tích hợp vào quá trình phát triển phần mềm.

Cụ thể, báo cáo tập trung vào ba vấn đề chính:
\begin{enumerate}
    \item Phân tích cách GenAI làm thay đổi trọng tâm của vòng đời phát triển phần mềm, đặc biệt là sự dịch chuyển từ các công việc mang tính lặp lại sang các hoạt động lập kế hoạch, thiết kế và ra quyết định kỹ thuật.
    \item Phân tích workflow cộng tác Người–AI gồm 6 bước được đề xuất trong nghiên cứu và vai trò của các nhóm kỹ năng khác nhau trong từng bước.
    \item Làm rõ vai trò của Human-in-the-loop, trong đó nhà phát triển không chỉ đưa yêu cầu cho AI mà còn liên tục đánh giá, hiệu chỉnh, cải thiện và chịu trách nhiệm đối với kết quả cuối cùng.
\end{enumerate}

Phạm vi của báo cáo tập trung vào góc nhìn của kỹ sư phần mềm trong quá trình sử dụng GenAI, thay vì đánh giá khả năng thay thế hoàn toàn con người của AI. Đây cũng phù hợp với kết luận của nghiên cứu rằng AI đang làm thay đổi và mở rộng phạm vi công việc của nhà phát triển theo hướng augmentation hơn là đơn thuần thay thế con người.


\newpage
\section{VÍ DỤ THỰC TẾ}
\input{2. VÍ DỤ THỰC TẾ/2.1 Tình huống từ tài liệu - Ca Alice (Task 2D)}
\input{2. VÍ DỤ THỰC TẾ/2.2 Tình huống mở rộng - Task 5B và 11D}
\input{2. VÍ DỤ THỰC TẾ/2.3 Tình huống minh họa ở sinh viên}
\input{2. VÍ DỤ THỰC TẾ/2.4 Đối chiếu các tình huống}

\newpage
\section{PHÂN TÍCH VẤN ĐỀ}
\subsection{Sự chuyển dịch trọng tâm trong SDLC}

Sự phát triển của Generative AI đang làm thay đổi cách thức thực hiện nhiều hoạt động trong vòng đời phát triển phần mềm (SDLC). Trước đây, kỹ sư phần mềm dành phần lớn thời gian trực tiếp thực hiện các tác vụ như viết mã nguồn, tìm lỗi, kiểm thử và xử lý các vấn đề kỹ thuật. Khi các công cụ AI có khả năng hỗ trợ tạo mã, giải thích mã, phát hiện lỗi và đề xuất giải pháp ngày càng phổ biến, một phần công việc thực thi có thể được AI hỗ trợ hoặc tự động hóa.

Tuy nhiên, điều này không đồng nghĩa với việc vai trò của kỹ sư phần mềm bị thay thế hoàn toàn. Trọng tâm công việc có xu hướng chuyển từ trực tiếp thực hiện từng tác vụ kỹ thuật sang xác định vấn đề, phân rã yêu cầu, hướng dẫn AI, đánh giá kết quả, kiểm thử/hiệu chỉnh, tích hợp và chịu trách nhiệm đối với sản phẩm cuối cùng. Kỹ sư cần có khả năng xác định khi nào nên sử dụng AI, cung cấp ngữ cảnh và yêu cầu phù hợp, đồng thời kiểm tra xem kết quả do AI tạo ra có đáp ứng đúng yêu cầu hay không.

Sự chuyển dịch này khiến các kỹ năng truyền thống vẫn giữ vai trò quan trọng. Kiến thức về lập trình, kiến trúc phần mềm, kiểm thử, bảo mật và quản lý yêu cầu giúp kỹ sư đánh giá được chất lượng của kết quả do AI tạo ra thay vì chỉ chấp nhận kết quả một cách máy móc. Đồng thời, các kỹ năng mới như làm việc với AI, phân rã vấn đề, kiểm chứng đầu ra và đánh giá rủi ro trở nên quan trọng hơn trong quy trình phát triển.

Có thể khái quát sự thay đổi này như sau:

\begin{itemize}
    \item \textbf{Workflow truyền thống:} Yêu cầu $\rightarrow$ Phân tích $\rightarrow$ Thiết kế $\rightarrow$ Lập trình $\rightarrow$ Kiểm thử $\rightarrow$ Sửa lỗi $\rightarrow$ Hoàn thiện.
    \item \textbf{Workflow có AI hỗ trợ:} Yêu cầu $\rightarrow$ Phân rã vấn đề $\rightarrow$ Giao nhiệm vụ cho AI $\rightarrow$ Đánh giá đầu ra $\rightarrow$ Kiểm thử/hiệu chỉnh $\rightarrow$ Tích hợp và chịu trách nhiệm.
\end{itemize}

Như vậy, AI chủ yếu làm thay đổi cách kỹ sư thực hiện công việc, trong khi con người vẫn giữ vai trò quan trọng trong việc xác định mục tiêu, đưa ra quyết định và kiểm soát chất lượng.

\subsection{Chu trình cộng tác Người-AI 6 bước và vai trò Human-in-the-loop}

Dựa trên cách tiếp cận về kỹ năng và workflow của kỹ sư phần mềm trong bối cảnh AI \cite{kam2025what}, quá trình phát triển có thể được mô hình hóa thành một chu trình cộng tác giữa con người và AI gồm 6 bước:

\begin{enumerate}
    \item \textbf{Bước 1 -- Xác định yêu cầu (Identify):} Con người xác định mục tiêu, yêu cầu chức năng, ràng buộc kỹ thuật và tiêu chí đầu ra. Đây là bước cần sự hiểu biết về bài toán và bối cảnh thực tế, những yếu tố mà AI không nên tự quyết định hoàn toàn.
    \item \textbf{Bước 2 -- Phân rã và giao nhiệm vụ (Engage):} Con người chia bài toán lớn thành các nhiệm vụ nhỏ hơn và cung cấp cho AI các yêu cầu, ngữ cảnh hoặc hướng dẫn phù hợp. Khả năng phân rã vấn đề giúp AI tạo ra kết quả có phạm vi và mục tiêu rõ ràng hơn.
    \item \textbf{Bước 3 -- AI tạo hoặc đề xuất giải pháp (Generate/Propose):} AI hỗ trợ thực hiện các tác vụ như sinh mã nguồn, đề xuất thuật toán, giải thích mã, tạo test case hoặc đưa ra phương án xử lý. Ở bước này, AI đóng vai trò là công cụ hỗ trợ thực thi thay vì là chủ thể đưa ra quyết định cuối cùng.
    \item \textbf{Bước 4 -- Con người kiểm tra và đánh giá (Evaluate):} Kết quả từ AI phải được con người xem xét về tính đúng đắn, khả năng đáp ứng yêu cầu, chất lượng mã nguồn, bảo mật và các rủi ro tiềm ẩn. Đây là một bước quan trọng vì đầu ra của AI có thể chứa lỗi hoặc không phù hợp với ngữ cảnh cụ thể.
    \item \textbf{Bước 5 -- Kiểm thử và hiệu chỉnh (Calibrate \& Tweak):} Kết quả được đưa qua quá trình kiểm thử. Nếu phát hiện vấn đề, con người có thể điều chỉnh yêu cầu, cung cấp thêm thông tin cho AI hoặc tự sửa đổi kết quả. Quá trình này có thể lặp lại nhiều lần cho đến khi sản phẩm đạt yêu cầu.
    \item \textbf{Bước 6 -- Tích hợp và chịu trách nhiệm (Finalize):} Kết quả đã được kiểm chứng được tích hợp vào hệ thống thực tế. Con người vẫn chịu trách nhiệm cuối cùng về quyết định đưa sản phẩm vào sử dụng, bao gồm chất lượng, bảo mật, tính phù hợp và các tác động có thể phát sinh.
\end{enumerate}

Chu trình trên thể hiện rõ vai trò của Human-in-the-loop (HITL). Con người không chỉ tham gia ở đầu vào bằng cách đưa yêu cầu cho AI mà còn xuất hiện xuyên suốt quá trình: xác định mục tiêu $\rightarrow$ hướng dẫn AI $\rightarrow$ đánh giá $\rightarrow$ kiểm thử $\rightarrow$ quyết định sử dụng. AI đảm nhiệm phần lớn các tác vụ có tính hỗ trợ và sinh nội dung, trong khi con người đảm nhiệm các hoạt động đòi hỏi hiểu biết về bối cảnh, phán đoán, kiểm chứng và trách nhiệm.

Do đó, trong workflow có GenAI, mối quan hệ phù hợp không phải là ``Con người $\rightarrow$ AI $\rightarrow$ Kết quả'', mà là một vòng lặp cộng tác:
\[
\text{Con người} \rightarrow \text{AI} \rightarrow \text{Kiểm tra} \rightarrow \text{Hiệu chỉnh} \rightarrow \text{AI} \rightarrow \text{Kiểm tra} \rightarrow \dots \rightarrow \text{Kết quả cuối cùng}
\]

Điểm quan trọng của mô hình này là AI không loại bỏ con người khỏi quy trình phát triển phần mềm mà làm thay đổi vị trí và trọng tâm tham gia của con người. Kỹ sư phần mềm ngày càng cần chuyển từ việc chỉ tập trung vào ``viết mã'' sang khả năng định hướng, đánh giá, kiểm chứng và chịu trách nhiệm đối với kết quả do AI hỗ trợ tạo ra.

\input{3. PHÂN TÍCH VẤN ĐỀ/3.3 Khung năng lực chữ T cho kỹ sư phần mềm}
\input{3. PHÂN TÍCH VẤN ĐỀ/3.4 Nguy cơ xói mòn kỹ năng và thách thức đối với kỹ sư trẻ}
\input{3. PHÂN TÍCH VẤN ĐỀ/3.5 Khoảng cách năng lực giữa mô hình AI-enhanced developer và sinh viên}

\newpage
\section{KẾT LUẬN}
\input{4. KẾT LUẬN/4.1 Tổng kết các luận điểm chính}
\input{4. KẾT LUẬN/4.2 Khuyến nghị lộ trình phát triển năng lực cho sinh viên CNTT}

% --- Tài liệu tham khảo ---
\newpage
\printbibliography[heading=bibintoc, title={TÀI LIỆU THAM KHẢO}]
% \bibliographystyle{IEEEtran}
% \bibliography{references}

\end{document}